% Generated draft: requires mathematical and source review.
\documentclass[11pt]{article}
\usepackage[margin=25mm]{geometry}
\usepackage{fontspec}
\setmainfont{Latin Modern Roman}
\usepackage{amsmath,amssymb,mathtools,unicode-math}
\setmathfont{Latin Modern Math}
\usepackage{enumitem,needspace,xurl,hyperref}
\setlength{\parindent}{0pt}
\setlength{\parskip}{5pt}
\setlength{\emergencystretch}{3em}
\newcommand{\N}{\mathbb{N}}
\newcommand{\Z}{\mathbb{Z}}
\newcommand{\Q}{\mathbb{Q}}
\newcommand{\R}{\mathbb{R}}
\newcommand{\C}{\mathbb{C}}
\newcommand{\F}{\mathbb{F}}
\newcommand{\A}{\mathbb{A}}
\newcommand{\PP}{\mathbb{P}}
\newcommand{\E}{\mathbb{E}}
\newcommand{\eps}{\varepsilon}
\newcommand{\dd}{\,\mathrm{d}}
\begin{document}
\section*{UnsolvedMath ID: 30004869}
\textbf{Source problem code:} catalogue entry
\subsection{Definitions and mathematical statement}
Let \(U\) be a finite collection of connected graphs. The supplied source asks about the graph profile \(G_U\) and the exponents of binomial inequalities valid on \(G_U\). Denote the intended set of exponents by \(E_U\).

For a specified real coordinate space \(\R^d\), a set \(C\subseteq\R^d\) is a polyhedral cone if there exist a nonnegative integer \(r\) and a matrix \(A\in\R^{r\times d}\) such that
\[
C=\{v\in\R^d:Av\geq 0\},
\]
where the inequality is componentwise. It is a rational polyhedral cone if such a matrix can be chosen in \(\Q^{r\times d}\).

\textbf{Source question, subject to the unresolved definitions below.} For every finite collection \(U\) of connected graphs, is \(E_U\) a polyhedral cone? If so, must \(E_U\) be a rational polyhedral cone in its intended exponent coordinates?

\textbf{Unresolved formulation.} The supplied record does not define \(G_U\), the class of graphs used to construct it, the normalization of its coordinates, or whether a closure is taken. It also does not specify the allowed binomial inequalities, their coefficients and exponent domains, or how their exponents are represented as vectors. In particular, it does not say whether an exponent vector records two monomial exponent lists or their difference. The ambient dimension and coordinate convention for \(E_U\) therefore remain unspecified. If zero coordinates occur in the profile, the admissibility and evaluation of monomials at those coordinates must also be specified. These omissions prevent an unconditional precise construction of \(E_U\) from this record; the question above preserves the supplied target without selecting those missing conventions.

The adjacent extracted discussion concerns sum-of-squares profiles and is truncated. It does not provide the missing definition of the true profile \(G_U\) or establish an identification of its valid binomial inequalities with those of a sum-of-squares profile.
\subsection{Short English statement}
The source asks whether, for every finite collection of connected graphs, the exponents of all binomial inequalities valid on its graph profile form a polyhedral cone, and whether that cone must be rational polyhedral. The supplied record omits the profile and exponent conventions needed to make this fully precise. Its open status is supplied dated curation.
\subsection{Sources}
\begin{itemize}
\item Local UnsolvedMath registry, supplied source record: permanent ID \(30004869\), title \emph{Polyhedral Cones of Graph Profile Inequality Exponents}, source code \texttt{catalogue entry}. The extracted question is labelled Question 2. No individual proposer is supplied; the proposed year is recorded as 2021.
\item Supplied authoritative source citation: \emph{Geometry and Optimization in Quantum Information (2022)}, Oberwolfach Reports, pp.\ 2665--2720, DOI \texttt{10.4171/owr/2021/49}. \url{https://doi.org/10.4171/owr/2021/49}.
\item Blekherman--Raymond--Singh--Thomas, supplied literature reference: \url{https://arxiv.org/abs/2004.05207}.
\item Published paper, supplied literature reference: \url{https://doi.org/10.1090/tran/8643}.
\item OWR report, also listed among the supplied literature sources: \url{https://doi.org/10.4171/owr/2021/49}.
\item Literature assessment dated exactly \texttt{2026-08-22}: the supplied curation labels the problem open and reports that Blekherman, Raymond, Singh, and Thomas prove convexity and compute important special cases while leaving the general polyhedral and rational-polyhedral questions open. This is supplied dated curation, not independently verified here; its time-sensitive status and cited claims require independent checking.
\item The registry labels the statement \texttt{corrected\_verified} and reports a direct check of the official report. That verification description is supplied provenance, not a source visit or independent verification performed for this draft.
\end{itemize}
\end{document}
